\section{Right-edge arithmetic source and exact source algebra}
\label{sec:right-edge}

\subsection{Hermitian reflection and elimination of the self block}
Differentiate the functional equation $\zeta(s)=\chi(s)\zeta(1-s)$. Since
$f(s)=-\frac12\chi'/\chi(s)$ and $f(s)=f(1-s)$,
\begin{equation}
\boxed{Z_1(s)=-\chi(s)Z_1(1-s).}
\label{eq:Z1-FE}
\end{equation}
and hence
\begin{equation}
\boxed{
L_1(s)+L_1(1-s)=\frac{\chi'}{\chi}(s),
\qquad L_1:=\frac{Z_1'}{Z_1}.
}
\label{eq:Z1-log-FE}
\end{equation}
The geometric reflection of the two vertical sides is
$\iota(s)=1-\bar s$.  Conjugation symmetry gives
\begin{equation}
\boxed{
-L_1(1-\bar s)
=\overline{L_1(s)-\frac{\chi'}{\chi}(s)}.
}
\label{eq:L1-hermitian-reflection}
\end{equation}
Put
\[
H(s):=f(s)+L_1(s)=\frac{\mathscr Z''}{\mathscr Z'}(s).
\]
The Hardy reflection law gives
\begin{equation}
\boxed{H(\iota(s))=-\overline{H(s)},}
\label{eq:hardy-source-adjoint}
\end{equation}
and likewise
\begin{equation}
\boxed{
C_{\mathscr Z}^{(p)}(\iota(s))
=\overline{C_{\mathscr Z}^{(p)}(s)},
\qquad
\mathcal B_{v,w}(\iota(s))
=\overline{\mathcal B_{w,v}(s)}.
}
\label{eq:curvature-packet-adjoint}
\end{equation}
If $\Gamma_c^+$ is the right side oriented upward and
$\Gamma_{1-c}^-$ the left side with boundary orientation, then the orientation
reversal cancels the minus sign in \eqref{eq:hardy-source-adjoint}. Therefore
\begin{equation}
\boxed{
\mathfrak I_{\rm left}(v,w)
=\overline{\mathfrak I_H(w,v)}
=\mathfrak I_H^*(v,w).
}
\label{eq:right-left-matrix-adjoint}
\end{equation}
Only after this matrix identity do we expand $H=f+L_1$.  Writing $I_f$ for the
right-edge form with scalar source $f$ gives
\begin{equation}
\boxed{
I_{\rm right}+I_{\rm right}^*+J_{\rm arch},
\qquad
J_{\rm arch}=I_f+I_f^*,
\qquad
2f=-\frac{\chi'}{\chi}.
}
\label{eq:J-I-source-split}
\end{equation}

\begin{lemma}[Holomorphic elimination of the archimedean self source]
\label{lem:holomorphic-self-elimination}
For every fixed normalized curvature derivative $0\le j\le2$,
\begin{equation}
\boxed{
\partial_p^j\oint_{\partial\mathcal R_T^{\rm sym}}
 f(s)C_{\mathscr Z}^{(p)}(s)\mathcal B_{v,w}(s)\,ds=0.
}
\label{eq:holomorphic-self-elimination}
\end{equation}
Consequently the complete $f$-self block is removed from the principal
arithmetic model before any AFE, stationary, packet, or HLZ reduction.  Its
vertical contribution is the negative of its horizontal contribution, and the
latter is covered on the retained packet space by
Lemma~\ref{lem:horizontal-decay}.
\end{lemma}

\begin{proof}
The factors $f$, $C_{\mathscr Z}^{(p)}$, and $\mathcal B_{v,w}$ are
holomorphic in the high rectangle, and the same is true after the fixed
$p$-derivatives. Cauchy's theorem gives
\eqref{eq:holomorphic-self-elimination}. Splitting the closed contour into
vertical and horizontal pieces gives the second assertion. No approximate
functional equation is involved.
\end{proof}

\subsection{Exact log-only right-edge split}
For bounded normalized curvature $p$, put $a=p/L$. On the right edge
\[
\mathscr Z(s+a)\mathscr Z(s-a)
=\Gamma_a(s)\zeta(s+a)\zeta(s-a),
\]
with
\[
\Gamma_0(s)=\chi(1-s),
\qquad
\partial_a\Gamma_a|_{a=0}=0,
\qquad
\partial_a^2\Gamma_a|_{a=0}=2f'(s)\chi(1-s).
\]
Since $f'(\sigma+it)=O(T^{-1})$ on fixed strips, every term containing this
second derivative belongs to the regular $f'$ class.

Fix
\begin{equation}
P(s):=-\frac{\zeta'}\zeta(s),
\qquad
Q(s):=-P'(s)=\sum_{n\ge1}\frac{\Lambda(n)\log n}{n^s}
\quad(\Re s>1).
\label{eq:PQ-convention}
\end{equation}
Then
\begin{equation}
\boxed{
L_1(s)=-P(s)+\frac{f'(s)+Q(s)}{f(s)-P(s)}.
}
\label{eq:L1-safe-line}
\end{equation}
Define
\begin{equation}
\boxed{H_{\log}(s):=-2P(s)+P(s-2/L),}
\label{eq:Hlog}
\end{equation}
and
\begin{equation}
\boxed{
\Delta H_{\rm dir}(s)
:=P(s)+\frac{Q(s)}{f(s)-P(s)}-P(s-2/L).
}
\label{eq:deltaH-dir}
\end{equation}
Then the exact identity is
\begin{equation}
\boxed{
L_1(s)=H_{\log}(s)+\Delta H_{\rm dir}(s)
      +\frac{f'(s)}{f(s)-P(s)}.
}
\label{eq:L1-logonly-split}
\end{equation}
Thus the principal arithmetic source contains neither an $f$-self term nor a
frozen block parameter.

\begin{lemma}[Exact curvature Dirichlet source]
\label{lem:curvature-source}
Put
\[
\mathcal C(s):=\zeta(s)^2Q(s).
\]
Then
\begin{equation}
\boxed{
C_{\mathscr Z}(s)
=q(s)^2\bigl(\mathcal C(s)+f'(s)\zeta(s)^2\bigr),
\qquad q(s)^2=\chi(1-s).
}
\label{eq:curvature-source}
\end{equation}
and
\begin{equation}
\mathcal C(s)=\sum_{n\ge1}\frac{c(n)}{n^s},
\qquad
\boxed{c(n)=\frac12\sum_{ab=n}\log^2(a/b)\ge0.}
\label{eq:curvature-coefficient}
\end{equation}
\end{lemma}

\begin{proof}
Since $\mathscr Z=q\zeta$ and $q'/q=f$,
\[
\mathscr Z''=q\bigl(\zeta''+2f\zeta'+(f'+f^2)\zeta\bigr),
\]
so
\[
C_{\mathscr Z}
=q^2\bigl(\zeta\zeta''-(\zeta')^2+f'\zeta^2\bigr).
\]
Because $(\zeta'/\zeta)'=Q$, the first two terms are $\zeta^2Q$.
The coefficient identity follows by symmetrizing
\[
\sum_{ab=n}\bigl((\log b)^2-\log a\log b\bigr).
\]
\end{proof}

\begin{lemma}[Product-level $f'$ separation]
\label{lem:logonly-fprime-separation}
After holomorphic self elimination, the right-edge product is exactly
\begin{equation}
\boxed{
q^2\left[\mathcal C H_{\log}
 +\mathcal C\Delta H_{\rm dir}
 +R_{f'}^{\log}\right],
}
\label{eq:logonly-product-split}
\end{equation}
where
\begin{equation}
\boxed{
R_{f'}^{\log}
=\frac{f'}{f-P}\mathcal C
 +f'\zeta^2H_{\log}
 +f'\zeta^2\Delta H_{\rm dir}
 +\frac{(f')^2}{f-P}\zeta^2.
}
\label{eq:logonly-fprime-source}
\end{equation}
Every term in $R_{f'}^{\log}$ contains at least one $f'$.
\end{lemma}

\begin{proof}
Multiply
$C_{\mathscr Z}/q^2=\mathcal C+f'\zeta^2$ by
\eqref{eq:L1-logonly-split}.  This gives the displayed identity exactly.
\end{proof}

\begin{lemma}[Second-order smallness of the direct residual]
\label{lem:deltaH-dir-small}
Fix $\Re s=1+\delta_0>1$. Uniformly for $t\asymp T$ and every fixed source
derivative,
\begin{equation}
\boxed{
\partial_s^j\Delta H_{\rm dir}(s)=O_{\delta_0,j}(L^{-2}),
\qquad j\ \text{fixed}.
}
\label{eq:deltaH-dir-small}
\end{equation}
The same holds after the fixed normalized curvature derivatives, up to a fixed
polylogarithmic factor.
\end{lemma}

\begin{proof}
On the safe line, $P,Q$ and their fixed derivatives are $O(1)$, whereas
$f=L/2+O(1)$. Thus
\[
\frac1{f-P}=\frac2L+O(L^{-2}).
\]
With $\delta=2/L$ and $P'=-Q$,
\[
P(s-\delta)=P(s)+\frac2LQ(s)+O(L^{-2}),
\qquad
\frac{Q(s)}{f(s)-P(s)}=\frac2LQ(s)+O(L^{-2}).
\]
Substitution proves the zeroth-order estimate. Fixed differentiation uses the
same safe-line bounds; after the leading constant $2/L$ is removed, every
derivative of $(f-P)^{-1}$ is $O(L^{-2})$ or smaller.
\end{proof}

The unique order-one source sent to the scalar HLZ master is therefore
\[
\boxed{\mathcal C(s)H_{\log}(s).}
\]
The two remaining families are $\mathcal C\Delta H_{\rm dir}$ and
$R_{f'}^{\log}$, both assigned directly to the relative-regular ledger.

\subsection{Legacy frozen-HLP comparison (not used in the main proof)}
\label{sec:legacy-frozen-hlp}
The following notation is retained only because several auxiliary estimates in
Section~\ref{sec:hlp-local} were originally formulated for the older frozen-HLP
comparison.  None of the objects in this subsection contributes to $P_T$ or to
the proof of Theorem~\ref{thm:main}.

Put
\begin{equation}
\boxed{
\widehat H_{\rm ex}(s)
=f(s)-P(s)+\frac{Q(s)}{f(s)-P(s)}.
}
\label{eq:Hex-variable}
\end{equation}
On a smooth height block define
\begin{equation}
\boxed{F_b:=f(1/2+i\tau_b)\in\mathbb R,}
\label{eq:Fb-real-freeze}
\end{equation}
\begin{equation}
\boxed{H_{\rm ex}(s;F_b)=F_b-P(s)+\frac{Q(s)}{F_b-P(s)},}
\label{eq:Hex}
\end{equation}
\begin{equation}
\boxed{H_{\rm mod}(s;F_b)=F_b-2P(s)+P(s-2/L),}
\label{eq:Hmod}
\end{equation}
and
\begin{equation}
\boxed{
\Delta H(s;F_b)
=P(s)+\frac{Q(s)}{F_b-P(s)}-P(s-2/L).
}
\label{eq:deltaH}
\end{equation}
For compatibility with the older auxiliary $f'$ bookkeeping, also set
\begin{equation}
\boxed{
q^2\left[\mathcal C\widehat H_{\rm ex}+R_{f'}^{\rm src}\right]
}
\label{eq:exact-fprime-separation}
\end{equation}
with
\begin{equation}
\boxed{
R_{f'}^{\rm src}
=f'\zeta^2\widehat H_{\rm ex}
+\frac{f'}{f-P}\mathcal C
+\frac{(f')^2}{f-P}\zeta^2.
}
\label{eq:fprime-source}
\end{equation}
These identities are auxiliary aliases only.

From this point, whenever the legacy HLP comparison is discussed, the spectral
variable $s$ and the independent Perron variable $z$ are kept distinct.

\begin{lemma}[Legacy locator for the translated model scale]
\label{lem:principal-pole-cluster}
Write $z=1+u$ and assume $F_b=L/2+O(1)$.  The frozen ratio
$H_{\rm ex}=F_b-P+Q/(F_b-P)$ has, as a meromorphic resummation near $u=0$,
a pole at $u=0$ of residue $-2$ and one zero-denominator pole
\[
u_*(F_b)=F_b^{-1}+O(F_b^{-2})=\frac2L+O(L^{-2})
\]
of residue $+1$.  The model $H_{\rm mod}$ has principal parts
$-2/u+1/(u-2/L)$.  This locator is not used by the revised principal route.
\end{lemma}

\begin{proof}
Use
\[
P(1+u)=u^{-1}+p_0+O(u),
\qquad
Q(1+u)=u^{-2}+O(1).
\]
For $D(u)=F_b-P(1+u)$, Rouch\'e's theorem on a circle of radius
$O(F_b^{-2})$ about $F_b^{-1}$ gives one zero
$u_*=F_b^{-1}+O(F_b^{-2})$. Since $D'=Q$, the residue of $Q/D$ there is one.
At $u=0$, both $-P$ and $Q/(F_b-P)$ contribute residue $-1$.
\end{proof}
